\section[Proof of Proposition 3.11]{Proof of \Cref{th:covariance_operator}}\label{sec:proof_Mercer_operator}

As $\mathcal{U}$ is compact and $k$ continuous, $k$ is bounded on $\mathcal{U}\times\mathcal{U}$, and $\mathcal{K}$ is compact.

Let $\phi,\phi'\in\mathcal{L}^2(\mathcal{U},\mathbb{R}^{d_z})$. Cauchy-Schwarz inequality ensures that there exists a constant $c\in\mathbb{R}_+$ such that:
\begin{equation}
	\iint_{\mathcal{U}\times\mathcal{U}}{\left|\phi(y)^\top k(x,y)^\top\phi'(x)\right|\:dx\:dy}\leq c\:\sqrt{\langle\phi|\phi\rangle}\:\sqrt{\langle\phi'|\phi'\rangle}<+\infty
\end{equation}

Fubini's integration theorem \cite{bib:Tonelli} can thus be used to give:
\begin{equation}
	\begin{aligned}
	\langle\mathcal{K}\phi|\phi'\rangle&=\int_{\mathcal{U}}{(\mathcal{K}\phi)(x)^\top\phi'(x)\:dx}\\
	&=\int_{\mathcal{U}}{\left(\int_{\mathcal{U}}{k(x,y)\:\phi(y)\:dy}\right)^\top\phi'(x)\:dx}\\
	&=\int_{\mathcal{U}}{\int_{\mathcal{U}}{\phi(y)^\top k(x,y)^\top\phi'(x)\:dy}\:dx}\\
	&=\int_{\mathcal{U}}{\int_{\mathcal{U}}{\phi(y)^\top k(y,x)\:\phi'(x)\:dx}\:dy}\\
	&=\int_{\mathcal{U}}{\phi(y)^\top\left(\int_{\mathcal{U}}{k(y,x)\:\phi'(x)\:dx}\right)dy}\\
	&=\int_{\mathcal{U}}{\phi(y)^\top(\mathcal{K}\phi')(y)\:dy}\\
	&=\langle\phi|\mathcal{K}\phi'\rangle
	\end{aligned}
\end{equation}

Let $\phi:\mathcal{U}\to\mathbb{R}^{d_z}$ be a continuous function, and let $n\in\mathbb{N}^*$. We define $\left(\mathcal{U}_i^{(n)}\right)_{i\in[\![1,n]\!]}$ a partition of $\mathcal{U}$ such that the diameter of each $\mathcal{U}_i^{(n)},\:i\in[\![1,n]\!]$ tends to $0$ when $n\to+\infty$. Here, the diameter of a subset of the Euclidean space $\mathbb{R}^{d_x}$ is defined as the maximum distance between two points of this subset. For $i\in[\![1,n]\!]$, we define $x_i^{(n)}\in\mathcal{U}_i^{(n)}$, and we let:
\begin{equation}
	a_i^{(n)}\triangleq\left(\int_{\mathcal{U}_i^{(n)}}{dx}\right)\phi\!\left(x_i^{(n)}\right)
\end{equation}

By continuity of $\phi$, the function:
\begin{equation}
	(x,y)\mapsto \phi(x)^\top k(x,y)\:\phi(y)
\end{equation}
is continuous on the compact $\mathcal{U}\times\mathcal{U}$, so its Riemann sums converge:
\begin{equation}
	\sum_{i=1}^n\sum_{j=1}^n\left(a_i^{(n)}\right)^\top k\left(x_i^{(n)},x_j^{(n)}\right)a_j^{(n)}\underset{n\to+\infty}{\longrightarrow}\iint_{\mathcal{U}\times\mathcal{U}}{\phi(x)^\top k(x,y)\:\phi(y)\:dx\:dy}=\langle\phi|\mathcal{K}\phi\rangle
\end{equation}

The SPSDness of $k$ ensures that each Riemann sum is non-negative. By passing to the limit, it follows that:
\begin{equation}
	\langle\phi|\mathcal{K}\phi\rangle\geq0\label{eq:fKf_geq0}
\end{equation}
for every continuous function $\phi$.

Finally, recall that the set of continuous functions of $\mathcal{U}$ is dense in $\mathcal{L}^2(\mathcal{U},\mathbb{R}^{d_z})$. By continuity of:
\begin{equation}
	\phi\in\mathcal{L}^2(\mathcal{U},\mathbb{R}^{d_z})\mapsto\langle\phi|\mathcal{K}\phi\rangle
\end{equation}
the inequality \eqref{eq:fKf_geq0} extends to all $\phi\in\mathcal{L}^2(\mathcal{U},\mathbb{R}^{d_z})$.

\section[Proof of Proposition 3.12]{Proof of \Cref{th:Mercer}}\label{sec:proof_Mercer}

By \Cref{th:covariance_operator}, $\mathcal{K}$ is a compact and self-adjoint operator on $\mathcal{L}^2(\mathcal{U},\mathbb{R}^{d_z})$. The spectral theorem \cite{bib:spectral-theorem} thus states that there exists an orthonormal basis $(\psi_i)_{i\in\mathbb{N}^*}$ of $\mathcal{L}^2(\mathcal{U},\mathbb{R}^{d_z})$ consisting of eigenfunctions of $\mathcal{K}$, with each corresponding eigenvalue $(\lambda_i)_{i\in\mathbb{N}^*}$ being real. Up to reordering, the sequence of eigenvalues can be considered to be non-increasing.

Since $k$ is uniformly continuous on the compact $\mathcal{U}\times\mathcal{U}$, $\mathcal{K}\phi$ is continuous for every $\phi\in\mathcal{L}^2(\mathcal{U},\mathbb{R}^{d_z})$. In particular, for each $i\in\mathbb{N}^*$ such that $\lambda_i\neq0$, the eigenfunction:
\begin{equation}
	\psi_i=\lambda_i^{-1}\mathcal{K}\psi_i
\end{equation}
is continuous. Without loss of generality, it must be noted that the $\psi_i$ associated with a null eigenvalue do not contribute to \eqref{eq:Mercer} and may thus be dismissed.

Since $(\psi_i)_{i\in\mathbb{N}^*}$ is an orthonormal basis of $\mathcal{L}^2(\mathcal{U},\mathbb{R}^{d_z})$, the family $(\psi_i\:\psi_j^\top)_{i,j\in\mathbb{N}^*}$ is an orthonormal basis of $\mathcal{F}$. Let $\langle\cdot|\cdot\rangle_{\mathcal{U}\times\mathcal{U}}$ denote:
\begin{equation}
	(\phi,\phi')\in\mathcal{F}^2\mapsto\langle\phi|\phi'\rangle_{\mathcal{U}\times\mathcal{U}}\triangleq\iint_{\mathcal{U}\times\mathcal{U}}{\operatorname{tr}\left(\phi(x,y)^\top\phi'(x,y)\right)\:dx\:dy}
\end{equation}

Let $i,j\in\mathbb{N}^*$. We have:
\begin{equation}
	\begin{aligned}
		\langle k|\psi_i\:\psi_j^\top\rangle_{\mathcal{U}\times\mathcal{U}}&=\iint_{\mathcal{U}\times\mathcal{U}}{\operatorname{tr}\left(k(x,y)^\top\psi_i(x)\:\psi_j(y)^\top\right)dx\:dy}\\
		&=\iint_{\mathcal{U}\times\mathcal{U}}{\psi_j(y)^\top k(x,y)^\top\psi_i(x)\:dx\:dy}\\
		&=\int_{\mathcal{U}}{\psi_j(y)^\top\left(\int_{\mathcal{U}}{k(y,x)\:\psi_i(x)\:dx}\right)dy}\\
		&=\int_{\mathcal{U}}{\psi_j(y)^\top\mathcal{K}\psi_i(y)\:dy}\\
		&=\lambda_i\int_{\mathcal{U}}{\psi_j(y)^\top\:\psi_i(y)\:dy}\\
		&=\lambda_i\:\delta_{i,j}
	\end{aligned}
\end{equation}
where:
\begin{equation}
	\delta_{i,j}\triangleq\left\{\begin{aligned}
	1&\quad\text{if }i=j\\0&\quad\text{if }i\neq j
	\end{aligned}\right.
\end{equation}
is the Kronecker symbol.

Because $k\in\mathcal{F}$, it follows from the completeness of $(\psi_i\:\psi_j^\top)_{i,j\in\mathbb{N}^*}$ that:
\begin{equation}
	k_{r_{\mathrm{KL}}}(x,x')\triangleq\sum_{i=1}^{r_{\mathrm{KL}}}\lambda_i\:\psi_i(x)\:\psi_i(x')^\top
\end{equation}
converges to $k$ in the $\mathcal{F}$ sense.

It must be noted that $\mathcal{L}^2$ convergence only identifies $k$ up to a Lebesgue-null subset of $\mathcal{U}\times\mathcal{U}$, and gives no control at any specific point of the domain. In particular, it is not enough to justify substituting $k_{r_{\mathrm{KL}}}$ for $k$ in kriging scenarios. However, since $k$ is continuous on $\mathcal{U}\times\mathcal{U}$, the convergence can be extended to absolute and uniform convergence on the entire domain \cite{bib:Karhunen-Loeve}. The argument applies Dini's theorem \cite{bib:Dini} to the non-decreasing sequence $\left(x\mapsto\operatorname{tr}\:k_{r_{\mathrm{KL}}}(x,x)\right)_{r_{\mathrm{KL}}\in\mathbb{N}^*}$, whose limit is continuous on the compact $\mathcal{U}$, and extends the result off the diagonal through the Cauchy-Schwarz inequality.

Since $\mathcal{K}$ is SPSD, every eigenvalue is non-negative. Assume that there exists a rank $r_0\in\mathbb{N}^*$ such that $\lambda_i=0$ as soon as $i>r_0$. Then the decomposition \eqref{eq:Mercer} would reduce to the finite sum:
\begin{equation}
	k=\sum_{i=1}^{r_0}\lambda_i\:\psi_i\:\psi_i^\top
\end{equation}
whose Gram matrices have rank at most $r_0$ and are thus singular at any $n$ distinct points with $nd_z>r_0$, contradicting the SPDness of $k$ by \Cref{th:kernel_gram_equivalence}. Consequently:
\begin{equation}
	\lambda_1\geq\lambda_2\geq\cdots>0
\end{equation}

\section[Proof of Proposition 3.14]{Proof of \Cref{th:reduced_rank_kriging}}\label{sec:proof_reduced_rank_kriging}

As \eqref{eq:reduced_rank_Gram} is a rank $r_{\mathrm{KL}}$ perturbation of $A\triangleq I_{\tilde{N}}\otimes\tilde{R}$, the Woodbury inversion formula (see \Cref{th:Woodbury}) can be used to reduce the $\tilde{N}d_z\times\tilde{N}d_z$ inversion of $G(\tilde{\mathbf{x}})$ to an $r_{\mathrm{KL}}\times r_{\mathrm{KL}}$ one:
\begin{equation}
	G(\tilde{\mathbf{x}})^{-1}\approx A^{-1}-A^{-1}\Psi(\tilde{\mathbf{x}})\:\tilde{G}(\tilde{\mathbf{x}})^{-1}\Psi(\tilde{\mathbf{x}})^\top A^{-1}
\end{equation}

Computing the centred SK estimator from \Cref{th:simple_kriging} using this approximation and the push-through identity (see \Cref{th:push_through}) yields:
\begin{equation}
	\begin{aligned}
	\hat{z}_{\mathrm{SK}}(x,\tilde{\mathbf{x}},\tilde{\mathbf{z}},0)&=k(x,\tilde{\mathbf{x}})\:G(\tilde{\mathbf{x}})^{-1}\tilde{\mathbf{z}}\\
	&\approx\Psi(x)\:\Lambda\:\Psi(\tilde{\mathbf{x}})^\top\left(I-A^{-1}\Psi(\tilde{\mathbf{x}})\:\tilde{G}(\tilde{\mathbf{x}})^{-1}\Psi(\tilde{\mathbf{x}})^\top\right)A^{-1}\tilde{\mathbf{z}}\\
	&=\Psi(x)\:\Lambda\left(I-\Psi(\tilde{\mathbf{x}})^\top A^{-1}\Psi(\tilde{\mathbf{x}})\:\tilde{G}(\tilde{\mathbf{x}})^{-1}\right)\Psi(\tilde{\mathbf{x}})^\top A^{-1}\tilde{\mathbf{z}}
	\end{aligned}\label{eq:z_s_up_to_tilde_G}
\end{equation}

Recall that, by definition of $\tilde{G}(\tilde{\mathbf{x}})$:
\begin{equation}
	\tilde{G}(\tilde{\mathbf{x}})-\Lambda^{-1}=\Psi(\tilde{\mathbf{x}})^\top A^{-1}\Psi(\tilde{\mathbf{x}})
\end{equation}
which gives:
\begin{equation}
	\begin{aligned}
	I-\Psi(\tilde{\mathbf{x}})^\top A^{-1}\Psi(\tilde{\mathbf{x}})\:\tilde{G}(\tilde{\mathbf{x}})^{-1}&=I-\left(\tilde{G}(\tilde{\mathbf{x}})-\Lambda^{-1}\right)\tilde{G}(\tilde{\mathbf{x}})^{-1}\\
	&=I-I+\Lambda^{-1}\tilde{G}(\tilde{\mathbf{x}})^{-1}\\
	&=\Lambda^{-1}\tilde{G}(\tilde{\mathbf{x}})^{-1}
	\end{aligned}\label{eq:reduced_rank_lemma}
\end{equation}

Plugging this result into \eqref{eq:z_s_up_to_tilde_G} and using $A^{-1}=I_{\tilde{N}}\otimes\tilde{R}^{-1}$ finally gives:
\begin{equation}
	\hat{z}_{\mathrm{SK}}(x,\tilde{\mathbf{x}},\tilde{\mathbf{z}},0)\approx\Psi(x)\:\tilde{G}(\tilde{\mathbf{x}})^{-1}\Psi(\tilde{\mathbf{x}})^\top A^{-1}\tilde{\mathbf{z}}
\end{equation}

For $R_{\mathrm{SK}}$, reusing:
\begin{equation}
	k(x,\tilde{\mathbf{x}})\:G(\tilde{\mathbf{x}})^{-1}\approx\Psi(x)\:\tilde{G}(\tilde{\mathbf{x}})^{-1}\Psi(\tilde{\mathbf{x}})^\top A^{-1}
\end{equation}
together with \eqref{eq:reduced_rank_lemma} yields:
\begin{equation}
	\begin{aligned}
	R_{\mathrm{SK}}&=k(x,x)+R-k(x,\tilde{\mathbf{x}})\:G(\tilde{\mathbf{x}})^{-1}k(\tilde{\mathbf{x}},x)\\
	&\approx\Psi(x)\:\Lambda\:\Psi(x)^\top+R-\Psi(x)\:\tilde{G}(\tilde{\mathbf{x}})^{-1}\Psi(\tilde{\mathbf{x}})^\top A^{-1}\times\Psi(\tilde{\mathbf{x}})\:\Lambda\:\Psi(x)^\top\\
	&=\Psi(x)\:\Lambda\:\Psi(x)^\top+R-\Psi(x)\:\tilde{G}(\tilde{\mathbf{x}})^{-1}\left(\tilde{G}(\tilde{\mathbf{x}})-\Lambda^{-1}\right)\Lambda\:\Psi(x)^\top\\
	&=\Psi(x)\:\Lambda\:\Psi(x)^\top+R-\Psi(x)\left(\Lambda-\tilde{G}(\tilde{\mathbf{x}})^{-1}\right)\Psi(x)^\top\\
	&=R+\Psi(x)\:\tilde{G}(\tilde{\mathbf{x}})^{-1}\Psi(x)^\top
	\end{aligned}
\end{equation}

Finally, when $\tilde{R}=\sigma^2I_{d_z}$ with $\sigma>0$, substituting $A^{-1}=\sigma^{-2}I_{\tilde{N}d_z}$ yields the announced simplifications. For such $\sigma>0$, the matrix:
\begin{equation}
	\Psi(\tilde{\mathbf{x}})^\top\Psi(\tilde{\mathbf{x}})+\sigma^2\Lambda^{-1}\succ0
\end{equation}
is SPD, as the sum of the SPSD matrix $\Psi(\tilde{\mathbf{x}})^\top\Psi(\tilde{\mathbf{x}})$ and the SPD matrix $\sigma^2\Lambda^{-1}$. The map:
\begin{equation}
	\sigma\in\mathbb{R}_+^*\mapsto\left(\Psi(\tilde{\mathbf{x}})^\top\Psi(\tilde{\mathbf{x}})+\sigma^2\Lambda^{-1}\right)^{-1}
\end{equation}
is thus well-defined and continuous, and extends to $\sigma=0$ when $\Psi(\tilde{\mathbf{x}})^\top\Psi(\tilde{\mathbf{x}})\succ0$, \textit{i.e.} when $\Psi(\tilde{\mathbf{x}})$ has full column rank. As $\hat{z}_{\mathrm{SK}}$ and $R_{\mathrm{SK}}$ depend on $\sigma$ only through this inverse and an explicit factor $\sigma^2$, both expressions remain valid at $\sigma=0$ under this condition.

\section[Proof of Proposition 3.15]{Proof of \Cref{th:KL}}\label{sec:proof_KL}

Let $i\in\mathbb{N}^*$. Since $\psi_i$ is bounded on the compact $\mathcal{U}$, Tonelli's theorem \cite{bib:Tonelli} and the continuity of $k$ give:
\begin{equation}
	\mathbb{E}\left[\int_{\mathcal{U}}{\|f(x)-m(x)\|^2\:dx}\right]=\int_{\mathcal{U}}{\operatorname{tr}k(x,x)\:dx}<\infty
\end{equation}
so that $f-m\in\mathcal{L}^2(\mathcal{U},\mathbb{R}^{d_z})$ a.s.. Each $\langle f-m|\psi_i\rangle$ is thus well-defined and finite. This integrability also justifies the exchanges of expectation and integration below.

Since any finite linear combination of $(\xi_i)_{i\in\mathbb{N}^*}$ is itself a linear combination of the values of $f$, they are jointly Gaussian. Let $j\in\mathbb{N}^*$. It comes that:
\begin{equation}
	\mathbb{E}[\xi_i]=\frac{1}{\sqrt{\lambda_i}}\int_{\mathcal{U}}{\mathbb{E}[f(x)-m(x)]^\top\psi_i(x)\:dx}=0
\end{equation}
and:
\begin{equation}
	\begin{aligned}
	\mathbb{E}[\xi_i\:\xi_j]&=\frac{1}{\sqrt{\lambda_i\:\lambda_j}}\iint_{\mathcal{U}\times\mathcal{U}}{\psi_i(x)^\top\mathbb{E}[(f(x)-m(x))(f(x')-m(x'))^\top]\:\psi_j(x')\:dx\:dx'}\\
	&=\frac{1}{\sqrt{\lambda_i\:\lambda_j}}\iint_{\mathcal{U}\times\mathcal{U}}{\psi_i(x)^\top k(x,x')\:\psi_j(x')\:dx\:dx'}\\
	&=\frac{1}{\sqrt{\lambda_i\:\lambda_j}}\int_{\mathcal{U}}{\psi_i(x)^\top\left(\int_{\mathcal{U}}{k(x,x')\:\psi_j(x')\:dx'}\right)\:dx}\\
	&=\frac{\lambda_j}{\sqrt{\lambda_i\:\lambda_j}}\int_{\mathcal{U}}{\psi_i(x)^\top\psi_j(x)\:dx}\\
	&=\frac{\lambda_j}{\sqrt{\lambda_i\:\lambda_j}}\langle\psi_i|\psi_j\rangle\\
	&=\delta_{i,j}
	\end{aligned}
\end{equation}
using the fact that $\psi_i$ is an eigenfunction of the Mercer operator $\mathcal{K}$. Together with the jointly Gaussian nature of $(\xi_i)_{i\in\mathbb{N}^*}$, this proves that they are i.i.d standard Gaussian variables.

Finally, for $r\in\mathbb{N}^*$, let:
\begin{equation}
	f_r:\left\{\begin{aligned}
	\mathcal{U}&\to\mathbb{R}^{d_z}\\
	x&\mapsto m(x)+\sum_{i=1}^r\sqrt{\lambda_i}\:\xi_i\:\psi_i(x)
	\end{aligned}\right.
\end{equation}

Fubini's theorem can be used to compute the following cross-covariance:
\begin{equation}
	\begin{aligned}
	\mathbb{E}\left[(f(x)-m(x))\:\xi_i\right]&=\frac{1}{\sqrt{\lambda_i}}\mathbb{E}[(f(x)-m(x))\langle f-m|\psi_i\rangle]\\
	&=\frac{1}{\sqrt{\lambda_i}}\mathbb{E}\left[(f(x)-m(x))\int_{\mathcal{U}}{(f(y)-m(y))^\top\psi_i(y)\:dy}\right]\\
	&=\frac{1}{\sqrt{\lambda_i}}\int_{\mathcal{U}}{\mathbb{E}\left[(f(x)-m(x))(f(y)-m(y))^\top\right]\psi_i(y)\:dy}\\
	&=\frac{1}{\sqrt{\lambda_i}}\int_{\mathcal{U}}{k(x,y)\:\psi_i(y)\:dy}\\
	&=\sqrt{\lambda_i}\:\psi_i(x)
	\end{aligned}
\end{equation}
which can be used to obtain:
\begin{equation}
	\begin{aligned}
	\mathbb{E}\left[\|f(x)-f_r(x)\|^2\right]&=\mathbb{E}\left[\|f(x)-m(x)\|^2\right]-2\sum_{i=1}^r\sqrt{\lambda_i}\:\psi_i(x)^\top\mathbb{E}\left[(f(x)-m(x))\:\xi_i\right]\\
	&\phantom{{}={}}+\sum_{i=1}^r\sum_{j=1}^r\sqrt{\lambda_i\:\lambda_j}\:\mathbb{E}[\xi_i\:\xi_j]\:\psi_i(x)^\top\psi_j(x)\\
	&=\operatorname{tr}k(x,x)-2\sum_{i=1}^r\lambda_i\operatorname{tr}\left(\psi_i(x)\:\psi_i(x)^\top\right)+\sum_{i=1}^r\lambda_i\operatorname{tr}\left(\psi_i(x)\:\psi_i(x)^\top\right)\hspace{-7pt}\\
	&=\operatorname{tr}k(x,x)-\sum_{i=1}^r\lambda_i\operatorname{tr}\left(\psi_i(x)\:\psi_i(x)^\top\right)\\
	&=\sum_{i=r+1}^{+\infty}\lambda_i\operatorname{tr}\left(\psi_i(x)\:\psi_i(x)^\top\right)
	\end{aligned}
\end{equation}

By Mercer's theorem, this last sum is the tail of:
\begin{equation}
	\operatorname{tr}k(x,x)=\sum_{i=1}^{+\infty}\lambda_i\:\operatorname{tr}\left(\psi_i(x)\:\psi_i(x)^\top\right)
\end{equation}
and therefore converges to 0. The partial sums thus converge to $f$ in mean square, uniformly over $\mathcal{U}$.

\section[Proof of Proposition 3.16]{Proof of \Cref{th:KL_converse}}\label{sec:proof_KL_converse}

For $r\in\mathbb{N}^*$, let:
\begin{equation}
	f_r:\left\{\begin{aligned}
	\mathcal{U}&\to\mathbb{R}^{d_z}\\
	x&\mapsto m(x)+\sum_{i=1}^r\sqrt{\lambda_i}\:\xi_i\:\psi_i(x)
	\end{aligned}\right.
\end{equation}
be the partial sum of the series.

Let $r'\in\mathbb{N}^*$ such that $r'>r$ and $x\in\mathcal{U}$. Since the $(\xi_i)_{i\in\mathbb{N}^*}$ are i.i.d standard Gaussian variables, it follows that:
\begin{equation}
	\begin{aligned}
	\mathbb{E}\left[\|f_r(x)-f_{r'}(x)\|^2\right]&=\sum_{i=r+1}^{r'}\lambda_i\:\|\psi_i(x)\|^2\\
	&=\sum_{i=r+1}^{r'}\lambda_i\:\operatorname{tr}\left(\psi_i(x)\:\psi_i(x)^\top\right)
	\end{aligned}
\end{equation}

As for \Cref{th:KL}, this sum is the tail of the convergent series:
\begin{equation}
	\operatorname{tr}k(x,x)=\sum_{i=1}^{+\infty}\lambda_i\:\operatorname{tr}\left(\psi_i(x)\:\psi_i(x)^\top\right)
\end{equation}
and tends to 0. The sequence $(f_r(x))_{r\in\mathbb{N}^*}$ is thus a Cauchy sequence in the space of square-integrable $\mathbb{R}^{d_z}$-valued random variables, which is complete. The sequence thus converges, and $f(x)$ is well-defined as its limit.

Let $n\in\mathbb{N}^*$, $x_1,\dots,x_n\in\mathcal{U}$ and $a_1,\dots,a_n\in\mathbb{R}^{d_z}$. For every $r\in\mathbb{N}^*$, the sum:
\begin{equation}
	\sum_{i=1}^na_i^\top f_r(x_i)
\end{equation}
is affine in $\xi_1,\dots,\xi_r$, and is thus Gaussian. Passing to the limit, this proves that:
\begin{equation}
	\sum_{i=1}^na_i^\top f(x_i)
\end{equation}
is also Gaussian, and so $f$ is a GP by \Cref{th:GP_definition}.

Since mean-square convergence carries the first two moments to the limit, the mean of $f$ is:
\begin{equation}
	\begin{aligned}
	\mathbb{E}[f(x)]&=\lim_{r\to+\infty}\mathbb{E}[f_r(x)]\\
	&=m(x)+\lim_{r\to+\infty}\sum_{i=1}^r\sqrt{\lambda_i}\:\mathbb{E}[\xi_i]\:\psi_i(x)\\
	&=m(x)
	\end{aligned}
\end{equation}
and its covariance is:
\begin{equation}
	\begin{aligned}
	\mathbb{E}\left[(f(x)-m(x))(f(x')-m(x'))^\top\right]&=\lim_{r\to+\infty}\sum_{i=1}^r\sum_{j=1}^r\sqrt{\lambda_i\:\lambda_j}\:\mathbb{E}[\xi_i\:\xi_j]\:\psi_i(x)\:\psi_j(x')^\top\\
	&=\lim_{r\to+\infty}\sum_{i=1}^r\sum_{j=1}^r\sqrt{\lambda_i\:\lambda_j}\:\delta_{i,j}\:\psi_i(x)\:\psi_j(x')^\top\\
	&=\sum_{i=1}^{+\infty}\lambda_i\:\psi_i(x)\:\psi_i(x')^\top\\
	&=k(x,x')
	\end{aligned}
\end{equation}

\section[Proof of Proposition 3.18]{Proof of \Cref{th:native_RKHS}}\label{sec:proof_native_RKHS}

\textbf{TODO: à relire}

\textbf{The form is well defined.} Let $\mathcal{H}_0$ denote the span of $\{k(x,\cdot),\:x\in\mathcal{U}\}$, and extend $\langle\cdot|\cdot\rangle_{\mathcal{H}_k}$ to it by bilinearity. For $\psi=\sum_{i=1}^na_i\:k(x_i,\cdot)$ and $\varphi=\sum_{j=1}^pb_j\:k(x'_j,\cdot)$ in $\mathcal{H}_0$, it comes that:
\begin{equation}
	\langle\psi|\varphi\rangle_{\mathcal{H}_k}=\sum_{i=1}^n\sum_{j=1}^pa_i\:b_j\:k(x_i,x'_j)=\sum_{j=1}^pb_j\:\psi(x'_j)=\sum_{i=1}^na_i\:\varphi(x_i)\label{eq:RKHS_bilinear}
\end{equation}
The last two expressions involve one of the two functions through its values alone, so neither representation matters and the form is well defined. It is symmetric and bilinear by construction, and positive since $k$ is SPSD. Taking $\varphi=k(x,\cdot)$ gives the reproducing property on $\mathcal{H}_0$:
\begin{equation}
	\forall\psi\in\mathcal{H}_0,\:\forall x\in\mathcal{U},\:\langle\psi|k(x,\cdot)\rangle_{\mathcal{H}_k}=\psi(x)\label{eq:reproducing_generators}
\end{equation}

\textbf{It is definite.} The Cauchy-Schwarz inequality holds for any positive semi-definite bilinear form, so \eqref{eq:reproducing_generators} bounds every evaluation:
\begin{equation}
	\forall\psi\in\mathcal{H}_0,\:\forall x\in\mathcal{U},\:|\psi(x)|^2\leq\langle\psi|\psi\rangle_{\mathcal{H}_k}\:k(x,x)\label{eq:RKHS_evaluation_bound}
\end{equation}
A function of null norm thus vanishes everywhere, and $\langle\cdot|\cdot\rangle_{\mathcal{H}_k}$ is an inner product on $\mathcal{H}_0$.

\textbf{The completion is a space of functions.} \Cref{eq:RKHS_evaluation_bound} makes every Cauchy sequence of $\mathcal{H}_0$ pointwise Cauchy, hence pointwise convergent, which sends each class of the abstract completion to a function on $\mathcal{U}$. That map is injective. Indeed, let $(\psi_n)_n$ be a Cauchy sequence converging pointwise to $0$, let $\varepsilon>0$, and let $N$ be such that $\|\psi_n-\psi_m\|_{\mathcal{H}_k}\leq\varepsilon$ for $n,m\geq N$. Fixing $n\geq N$ and writing $\psi_n=\sum_{i=1}^{n'}a_i\:k(x_i,\cdot)$, \eqref{eq:RKHS_bilinear} gives:
\begin{equation}
	\|\psi_n\|_{\mathcal{H}_k}^2=\langle\psi_n|\psi_n-\psi_m\rangle_{\mathcal{H}_k}+\sum_{i=1}^{n'}a_i\:\psi_m(x_i)
\end{equation}
whose first term is at most $\varepsilon\sup_n\|\psi_n\|_{\mathcal{H}_k}$, finite since a Cauchy sequence is bounded, and whose second is a finite sum tending to $0$ with $m$. Hence $\|\psi_n\|_{\mathcal{H}_k}\to0$, and only the zero class is sent to the zero function. Passing \eqref{eq:reproducing_generators} to the limit, $\mathcal{H}_k$ is an RKHS of reproducing kernel $k$.

\textbf{Uniqueness.} Let $\mathcal{H}$ be an RKHS of reproducing kernel $k$. It contains every $k(x,\cdot)$, hence $\mathcal{H}_0$, and its inner product agrees with $\langle\cdot|\cdot\rangle_{\mathcal{H}_k}$ on the generators, hence on $\mathcal{H}_0$ by bilinearity. A function of $\mathcal{H}$ orthogonal to $\mathcal{H}_0$ verifies $\psi(x)=\langle\psi|k(x,\cdot)\rangle_{\mathcal{H}}=0$ for every $x\in\mathcal{U}$, so $\mathcal{H}_0$ is dense in $\mathcal{H}$, which is therefore its completion $\mathcal{H}_k$.

\section[Proof of Proposition 4.3]{Proof of \Cref{th:errors_inverse_Gram}}\label{sec:proof_errors_inverse_Gram}

As established in \cref{sec:proof_UK}, the weight matrix $W$ of \eqref{eq:kriging_f} and the Lagrange multiplier $\xi$ of the unbiasedness constraints verify:
\begin{equation}
	\mathcal{G}(\tilde{\mathbf{x}})\begin{bmatrix}W\\\xi\end{bmatrix}=\bar{k}(\tilde{\mathbf{x}},x)
\end{equation}
whose first block row is $G(\tilde{\mathbf{x}})\:W+g(\tilde{\mathbf{x}})\:\xi=k(\tilde{\mathbf{x}},x)$. Let $e_{\tilde{\mathbf{x}}}$ and $e_{\tilde{\mathbf{x}}^r}$ denote the errors of the full and subsampled estimators:
\begin{equation}
	\left\{\begin{aligned}
	e_{\tilde{\mathbf{x}}}&\triangleq z-\hat{z}_{\mathrm{UK}}(x,\tilde{\mathbf{x}},\tilde{\mathbf{z}},g)\\
	e_{\tilde{\mathbf{x}}^r}&\triangleq z-\hat{z}_{\mathrm{UK}}(x,\tilde{\mathbf{x}}^r,\tilde{\mathbf{z}}^r,g)
	\end{aligned}\right.
\end{equation}

Since $e_{\tilde{\mathbf{x}}}$ is centred whatever the value of $\beta$, and $\varepsilon$ is independent of $\tilde{\mathbf{z}}$, it comes that:
\begin{equation}
	\mathbb{E}\left[\tilde{\mathbf{z}}\:e_{\tilde{\mathbf{x}}}^\top\right]=k(\tilde{\mathbf{x}},x)-G(\tilde{\mathbf{x}})\:W=g(\tilde{\mathbf{x}})\:\xi
\end{equation}

Therefore, $e_{\tilde{\mathbf{x}}}$ is uncorrelated with every linear combination $V^\top\tilde{\mathbf{z}}$ such that $V^\top g(\tilde{\mathbf{x}})=0$. The difference between the two interpolators can be decomposed as follows:
\begin{equation}
	\hat{z}_{\mathrm{UK}}(x,\tilde{\mathbf{x}},\tilde{\mathbf{z}},g)-\hat{z}_{\mathrm{UK}}(x,\tilde{\mathbf{x}}^r,\tilde{\mathbf{z}}^r,g)=e_{\tilde{\mathbf{x}}^r}-e_{\tilde{\mathbf{x}}}
\end{equation}

Thus:
\begin{equation}
	\begin{aligned}
	E_r(x)&=\mathbb{E}\left[\left(e_{\tilde{\mathbf{x}}^r}-e_{\tilde{\mathbf{x}}}\right)\left(e_{\tilde{\mathbf{x}}^r}-e_{\tilde{\mathbf{x}}}\right)^\top\right]\\
	&=\mathbb{E}\left[e_{\tilde{\mathbf{x}}^r}e_{\tilde{\mathbf{x}}^r}^\top\right]+\mathbb{E}\left[e_{\tilde{\mathbf{x}}}e_{\tilde{\mathbf{x}}}^\top\right]-\mathbb{E}\left[e_{\tilde{\mathbf{x}}}e_{\tilde{\mathbf{x}}^r}^\top\right]-\mathbb{E}\left[e_{\tilde{\mathbf{x}}^r}e_{\tilde{\mathbf{x}}}^\top\right]
	\end{aligned}
\end{equation}

Since both interpolators are unbiased, their difference is such a linear combination, and so:
\begin{equation}
	\mathbb{E}\left[e_{\tilde{\mathbf{x}}}\left(e_{\tilde{\mathbf{x}}^r}-e_{\tilde{\mathbf{x}}}\right)^\top\right]=0
\end{equation}

The last two terms are thus both equal to $\mathbb{E}\left[e_{\tilde{\mathbf{x}}}e_{\tilde{\mathbf{x}}}^\top\right]$. Moreover, the second moments of the errors of unbiased interpolators are their covariances. Therefore:
\begin{equation}
	E_r(x)=\mathbb{E}\left[e_{\tilde{\mathbf{x}}^r}e_{\tilde{\mathbf{x}}^r}^\top\right]-\mathbb{E}\left[e_{\tilde{\mathbf{x}}}e_{\tilde{\mathbf{x}}}^\top\right]=R_{\mathrm{UK}}(x,\tilde{\mathbf{x}}^r,g)-R_{\mathrm{UK}}(x,\tilde{\mathbf{x}},g)
\end{equation}

Using \Cref{th:subsampled_kriging} yields:
\begin{equation}
	\begin{aligned}
	E_r(x)&=\bar{k}(x,\tilde{\mathbf{x}})\:\mathcal{G}(\tilde{\mathbf{x}})^{-1}\bar{k}(\tilde{\mathbf{x}},x)-\bar{k}(x,\tilde{\mathbf{x}})\:\mathcal{G}_h(\tilde{\mathbf{x}}^r)^\dagger\:\bar{k}(\tilde{\mathbf{x}},x)\\
	&=\bar{k}(x,\tilde{\mathbf{x}})\:\mathfrak{D}_G\bar{k}(\tilde{\mathbf{x}},x)
	\end{aligned}
\end{equation}

For $e_f(x)$, direct calculations give:
\begin{equation}
	\begin{aligned}
	e_f(x)&=\left\|\hat{z}_{\mathrm{UK}}(x,\tilde{\mathbf{x}},\tilde{\mathbf{z}},g)-\hat{z}_{\mathrm{UK}}(x,\tilde{\mathbf{x}}^r,\tilde{\mathbf{z}}^r,g)\right\|_2\\
	&=\left\|\bar{k}(x,\tilde{\mathbf{x}})\:\mathcal{G}(\tilde{\mathbf{x}})^{-1}\bar{\mathbf{z}}-\bar{k}(x,\tilde{\mathbf{x}})\:\mathcal{G}_h(\tilde{\mathbf{x}}^r)^\dagger\:\bar{\mathbf{z}}\right\|_2\\
	&=\left\|\bar{k}(x,\tilde{\mathbf{x}})\:\mathfrak{D}_G\bar{\mathbf{z}}\right\|_2
	\end{aligned}
\end{equation}

\section[Proof of Proposition 4.6]{Proof of \Cref{th:explicit_residuals}}\label{sec:proof_explicit_residuals}

For the sake of brevity, we write $G_r$, $g_r$ and $M_r$ for $G(\tilde{\mathbf{x}}^r)$, $g(\tilde{\mathbf{x}}^r)$ and $M(\tilde{\mathbf{x}}^r)$. By \Cref{th:Schur}, the inverse of $A=\mathcal{G}(\tilde{\mathbf{x}}^r)$ is:
\begin{equation}
	\mathcal{G}(\tilde{\mathbf{x}}^r)^{-1}=\begin{bmatrix}G_r^{-1}-G_r^{-1}g_r\:M_r\:g_r^\top G_r^{-1}&G_r^{-1}g_r\:M_r\\M_r\:g_r^\top G_r^{-1}&-M_r\end{bmatrix}
\end{equation}

Let $x'\in\mathcal{U}$, and let $y$ and $t$ be two matrices with respectively $\tilde{N}^rd_z$ and $r_{\mathrm{UK}}d_z$ rows. It comes that:
\begin{equation}
	\bar{k}(x',\tilde{\mathbf{x}}^r)\:\mathcal{G}(\tilde{\mathbf{x}}^r)^{-1}\begin{bmatrix}y\\t\end{bmatrix}=k(x',\tilde{\mathbf{x}}^r)\:G_r^{-1}y-U(x')\:M_r\left(t-g_r^\top G_r^{-1}y\right)\label{eq:bordered_product}
\end{equation}

The three expressions of \eqref{eq:explicit_residuals} are obtained by applying \eqref{eq:bordered_product} to every element $x'$ of $\tilde{\mathbf{x}}^d$:
\begin{itemize}
	\item Taking $y=k(\tilde{\mathbf{x}}^r,x)$ and $t=g(x)^\top$ gives the expression of $k_{d|r}(x)$, since $t-g_r^\top G_r^{-1}y$ is then equal to $U(x)^\top$.
	\item Taking the columns of $k(\tilde{\mathbf{x}}^r,\tilde{\mathbf{x}}^d)$ and of $g(\tilde{\mathbf{x}}^d)^\top$ for $y$ and $t$ gives the expression of $S_{d|r}=D-B^\top A^{-1}B$ in the same way.
	\item Taking $y=\tilde{\mathbf{z}}^r$ and $t=0$ gives the expression of $\tilde{\mathbf{z}}_{d|r}$, since $g_r^\top G_r^{-1}\tilde{\mathbf{z}}^r=M_r^{-1}\hat{m}_{\mathrm{UK}}(\tilde{\mathbf{x}}^r,\tilde{\mathbf{z}}^r,g)$.
\end{itemize}

Similarly, $S^0_{d|r}$ is the Schur complement of $G(\tilde{\mathbf{x}})$ with respect to $G_r$, and \Cref{th:Schur} gives:
\begin{equation}
	G(\tilde{\mathbf{x}})^{-1}=\begin{bmatrix}G_r^{-1}&0\\0&0\end{bmatrix}+\begin{bmatrix}-G_r^{-1}k(\tilde{\mathbf{x}}^r,\tilde{\mathbf{x}}^d)\\I_{\tilde{N}^dd_z}\end{bmatrix}\left(S^0_{d|r}\right)^{-1}\begin{bmatrix}-G_r^{-1}k(\tilde{\mathbf{x}}^r,\tilde{\mathbf{x}}^d)\\I_{\tilde{N}^dd_z}\end{bmatrix}^\top
\end{equation}

Multiplying this expression by $g(\tilde{\mathbf{x}})^\top$ on the left and by $g(\tilde{\mathbf{x}})$ on the right yields:
\begin{equation}
	g(\tilde{\mathbf{x}})^\top G(\tilde{\mathbf{x}})^{-1}g(\tilde{\mathbf{x}})=g_r^\top G_r^{-1}g_r+U(\tilde{\mathbf{x}}^d)^\top\left(S^0_{d|r}\right)^{-1}U(\tilde{\mathbf{x}}^d)
\end{equation}
which is exactly \eqref{eq:trend_information}.

Finally, $S^0_{d|r}$ is SPD as the Schur complement of the SPD matrix $G(\tilde{\mathbf{x}})$ with respect to $G_r$ \cite{bib:matrix-analysis}. Since $U(\tilde{\mathbf{x}}^d)\:M_r\:U(\tilde{\mathbf{x}}^d)^\top\succeq0$, $S_{d|r}$ is thus SPD as well. Moreover, adding $g(\tilde{\mathbf{x}})\:\beta_0$ to the observations adds $U(\tilde{\mathbf{x}}^d)\:\beta_0$ to $\tilde{\mathbf{z}}^0_{d|r}$ and $\beta_0$ to $\hat{m}_{\mathrm{UK}}(\tilde{\mathbf{x}}^r,\tilde{\mathbf{z}}^r,g)$, which leaves $\tilde{\mathbf{z}}_{d|r}$ unchanged.

\section[Proof of Propositions 4.10 and 4.11]{Proof of \Cref{th:split_bound,th:computable_bound}}\label{sec:proof_split_bound}

For the sake of brevity, we write $M$, $M_r$ and $U_d$ for $M(\tilde{\mathbf{x}})$, $M(\tilde{\mathbf{x}}^r)$ and $U(\tilde{\mathbf{x}}^d)$. By \Cref{th:explicit_residuals}, we have:
\begin{equation}
	\left\{\begin{aligned}
	k_{d|r}(x)&=k^0_{d|r}(x)+U_d\:M_r\:U(x)^\top\\
	S_{d|r}&=S^0_{d|r}+U_d\:M_r\:U_d^\top\succeq S^0_{d|r}
	\end{aligned}\right.
\end{equation}
and thus $\|\cdot\|_{S_{d|r}}\leq\|\cdot\|_{S^0_{d|r}}$. Applying the triangle inequality to \Cref{th:errors_euclidean} gives:
\begin{equation}
	e_r(x)\leq\|k^0_{d|r}(x)\|_{S^0_{d|r}}+\left\|U_d\:M_r\:U(x)^\top\right\|_{S_{d|r}}\label{eq:split_triangle}
\end{equation}
with equality when $k^0_{d|r}(x)=0$.

By \eqref{eq:trend_information}, the Woodbury formula (see \Cref{th:Woodbury}) applied to $S_{d|r}$ yields:
\begin{equation}
	S_{d|r}^{-1}=\left(S^0_{d|r}\right)^{-1}-\left(S^0_{d|r}\right)^{-1}U_d\:M\:U_d^\top\left(S^0_{d|r}\right)^{-1}
\end{equation}
and thus, using \eqref{eq:trend_information} once more:
\begin{equation}
	\begin{aligned}
	U_d^\top S_{d|r}^{-1}U_d&=\left(M^{-1}-M_r^{-1}\right)-\left(M^{-1}-M_r^{-1}\right)M\left(M^{-1}-M_r^{-1}\right)\\
	&=\left(M^{-1}-M_r^{-1}\right)M\:M_r^{-1}\\
	&=M_r^{-1}-M_r^{-1}M\:M_r^{-1}
	\end{aligned}
\end{equation}

The second term of \eqref{eq:split_triangle} is therefore:
\begin{equation}
	\begin{aligned}
	\left\|U_d\:M_r\:U(x)^\top\right\|_{S_{d|r}}^2&=\left\|U(x)\:M_r\:U_d^\top S_{d|r}^{-1}U_d\:M_r\:U(x)^\top\right\|_{\mathrm{op}}\\
	&=\left\|U(x)\left(M_r-M\right)U(x)^\top\right\|_{\mathrm{op}}
	\end{aligned}
\end{equation}
which proves \Cref{th:split_bound}.

For \Cref{th:computable_bound}, it must be noted that $S^0_{d|r}-I_{\tilde{N}^d}\otimes\tilde{R}$ is the Schur complement of:
\begin{equation}
	k(\tilde{\mathbf{x}},\tilde{\mathbf{x}})+\begin{bmatrix}I_{\tilde{N}^r}&0\\0&0\end{bmatrix}\otimes\tilde{R}\succeq0
\end{equation}
with respect to $G(\tilde{\mathbf{x}}^r)$. Since this matrix is SPSD, so is its Schur complement \cite{bib:matrix-analysis}, and thus:
\begin{equation}
	S^0_{d|r}\succeq I_{\tilde{N}^d}\otimes\tilde{R}\succeq\tilde{\lambda}\:I_{\tilde{N}^dd_z}
\end{equation}

Consequently, we have:
\begin{equation}
	\left\{\begin{aligned}
	\|k^0_{d|r}(x)\|_{S^0_{d|r}}&\leq\frac{\|k^0_{d|r}(x)\|_{\mathrm{op}}}{\sqrt{\tilde{\lambda}}}\\
	M^{-1}&\preceq M_r^{-1}+\frac{U_d^\top U_d}{\tilde{\lambda}}
	\end{aligned}\right.
\end{equation}
where the second inequality follows from \eqref{eq:trend_information}. Since the inversion reverses the order of SPD matrices, applying these inequalities to \Cref{th:split_bound} concludes the proof.

\section[Proof of Proposition 4.13]{Proof of \Cref{th:entrywise_residual}}\label{sec:proof_entrywise}

Let $N(K)$ denote the matrix of the operator norms of the blocks of a matrix $K$, and let $v$ be a vector whose blocks have Euclidean norms gathered in $u$. Since $\|Kv\|_2\leq\|N(K)\:u\|_2$, it comes that:
\begin{equation}
	\|K\|_{\mathrm{op}}\leq\|N(K)\|_{\mathrm{op}}\leq\sqrt{\|N(K)\|_1\:\|N(K)\|_\infty}
\end{equation}
where $\|N(K)\|_1$ and $\|N(K)\|_\infty$ are the largest column and row sums of $N(K)$ \cite{bib:matrix-analysis}. Applied to the symmetric matrix $G(\tilde{\mathbf{x}}^j)$, this inequality gives the blockwise Gershgorin bound:
\begin{equation}
	\lambda_{\max}\left(G(\tilde{\mathbf{x}}^j)\right)\leq\lambda_j^+
\end{equation}

Similarly, applying it to \eqref{eq:entrywise_assumptions} gives:
\begin{equation}
	\left\{\begin{aligned}
	\|k(\tilde{\mathbf{x}}^d,\tilde{\mathbf{x}}^r)\|_{\mathrm{op}}&\leq\sqrt{\tilde{N}^r\tilde{N}^d}\:\alpha\\
	\|k(\tilde{\mathbf{x}}^r,x)\|_{\mathrm{op}}&\leq\sqrt{\tilde{N}^r}\:\kappa\\
	\|k(\tilde{\mathbf{x}}^d,x)\|_{\mathrm{op}}&\leq\sqrt{\tilde{N}^d}\:\eta\\
	\|g(\tilde{\mathbf{x}}^j)\|_{\mathrm{op}}&\leq\sqrt{\tilde{N}^j}\:g_{\max}\\
	\|\tilde{\mathbf{z}}^j-g(\tilde{\mathbf{x}}^j)\:\beta_0\|_2&\leq\sqrt{\tilde{N}^j}\:\tilde{z}_{\max}
	\end{aligned}\right.\label{eq:entrywise_norms}
\end{equation}

Since $G(\tilde{\mathbf{x}}^r)$ is SPD, it follows that:
\begin{equation}
	\begin{aligned}
	\|k^0_{d|r}(x)\|_{\mathrm{op}}&\leq\|k(\tilde{\mathbf{x}}^d,x)\|_{\mathrm{op}}+\frac{\|k(\tilde{\mathbf{x}}^d,\tilde{\mathbf{x}}^r)\|_{\mathrm{op}}\:\|k(\tilde{\mathbf{x}}^r,x)\|_{\mathrm{op}}}{\lambda_r^-}\\
	&\leq\sqrt{\tilde{N}^d}\:\frac{\lambda_r^-\eta+\alpha\tilde{N}^r\kappa}{\lambda_r^-}
	\end{aligned}
\end{equation}

The same steps, applied to $U(\tilde{\mathbf{x}}^d)$ and $U(x)$, give:
\begin{equation}
	\left\{\begin{aligned}
	\|U(\tilde{\mathbf{x}}^d)\|_{\mathrm{op}}&\leq\sqrt{\tilde{N}^d}\:g_{\max}\:\rho_r\\
	\|U(x)\|_{\mathrm{op}}&\leq g_{\max}\:\frac{\lambda_r^-+\tilde{N}^r\kappa}{\lambda_r^-}
	\end{aligned}\right.
\end{equation}

Furthermore, for every unit vector $v\in\mathbb{R}^{r_{\mathrm{UK}}d_z}$, the Rayleigh quotient \cite{bib:matrix-analysis} and the definitions of $\lambda_r^+$ and $\lambda_r^g$ give:
\begin{equation}
	v^\top M(\tilde{\mathbf{x}}^r)^{-1}v=\left(g(\tilde{\mathbf{x}}^r)\:v\right)^\top G(\tilde{\mathbf{x}}^r)^{-1}\left(g(\tilde{\mathbf{x}}^r)\:v\right)\geq\frac{\|g(\tilde{\mathbf{x}}^r)\:v\|_2^2}{\lambda_r^+}\geq\frac{\lambda_r^g}{\lambda_r^+}
\end{equation}
and thus $\|M(\tilde{\mathbf{x}}^r)\|_{\mathrm{op}}\leq\lambda_r^+/\lambda_r^g$. Applying the triangle inequality to the expression of $k_{d|r}(x)$ given by \Cref{th:explicit_residuals}, and combining these inequalities, yields the bound on $\|k_{d|r}(x)\|_{\mathrm{op}}$.

Since $S_{d|r}\succeq S^0_{d|r}$ by \Cref{th:explicit_residuals}, we have $\lambda_{\min}(S_{d|r})\geq\lambda_{\min}(S^0_{d|r})$. Moreover, for every unit vector $v\in\mathbb{R}^{\tilde{N}^dd_z}$:
\begin{equation}
	\begin{aligned}
	v^\top S^0_{d|r}\:v&=v^\top G(\tilde{\mathbf{x}}^d)\:v-\left(k(\tilde{\mathbf{x}}^r,\tilde{\mathbf{x}}^d)\:v\right)^\top G(\tilde{\mathbf{x}}^r)^{-1}\left(k(\tilde{\mathbf{x}}^r,\tilde{\mathbf{x}}^d)\:v\right)\\
	&\geq\lambda_d^--\frac{\|k(\tilde{\mathbf{x}}^r,\tilde{\mathbf{x}}^d)\|_{\mathrm{op}}^2}{\lambda_r^-}\\
	&\geq\frac{\lambda_r^-\lambda_d^--\alpha^2\tilde{N}^r\tilde{N}^d}{\lambda_r^-}
	\end{aligned}
\end{equation}

Since this lower bound does not depend on $v$, it also bounds $\lambda_{\min}(S^0_{d|r})$. It is positive by assumption, and so:
\begin{equation}
	\|S_{d|r}^{-1}\|_{\mathrm{op}}=\frac{1}{\lambda_{\min}(S_{d|r})}\leq s_{\max}
\end{equation}

Finally, since $\tilde{\mathbf{z}}_{d|r}$ is unchanged when a trend is added to the observations, \Cref{th:explicit_residuals} can be applied to the observations $\tilde{\mathbf{z}}-g(\tilde{\mathbf{x}})\:\beta_0$:
\begin{equation}
	\begin{aligned}
	\tilde{\mathbf{z}}_{d|r}&=\left(\tilde{\mathbf{z}}^d-g(\tilde{\mathbf{x}}^d)\:\beta_0\right)-k(\tilde{\mathbf{x}}^d,\tilde{\mathbf{x}}^r)\:G(\tilde{\mathbf{x}}^r)^{-1}\left(\tilde{\mathbf{z}}^r-g(\tilde{\mathbf{x}}^r)\:\beta_0\right)\\
	&\qquad-U(\tilde{\mathbf{x}}^d)\left(\hat{m}_{\mathrm{UK}}(\tilde{\mathbf{x}}^r,\tilde{\mathbf{z}}^r,g)-\beta_0\right)
	\end{aligned}
\end{equation}

The steps used for $k^0_{d|r}(x)$ bound the norm of the first two terms by $\sqrt{\tilde{N}^d}\:\rho_r\:\tilde{z}_{\max}$. Moreover, \eqref{eq:entrywise_norms} gives:
\begin{equation}
	\begin{aligned}
	\left\|\hat{m}_{\mathrm{UK}}(\tilde{\mathbf{x}}^r,\tilde{\mathbf{z}}^r,g)-\beta_0\right\|_2&=\left\|M(\tilde{\mathbf{x}}^r)\:g(\tilde{\mathbf{x}}^r)^\top G(\tilde{\mathbf{x}}^r)^{-1}\left(\tilde{\mathbf{z}}^r-g(\tilde{\mathbf{x}}^r)\:\beta_0\right)\right\|_2\\
	&\leq\frac{\lambda_r^+}{\lambda_r^g}\:\frac{\sqrt{\tilde{N}^r}\:g_{\max}\:\sqrt{\tilde{N}^r}\:\tilde{z}_{\max}}{\lambda_r^-}\\
	&=\frac{\tilde{N}^rg_{\max}\:\lambda_r^+}{\lambda_r^g\:\lambda_r^-}\:\tilde{z}_{\max}
	\end{aligned}
\end{equation}
which concludes the proof by the triangle inequality.
